\documentclass[11pt,letterpaper]{article}
\usepackage[T1]{fontenc}
\usepackage[utf8]{inputenc}
\usepackage[scaled=0.95]{helvet}\renewcommand{\familydefault}{\sfdefault}
\usepackage[margin=2cm]{geometry}
\usepackage{graphicx,xcolor,enumitem,booktabs,amsmath,siunitx}
\usepackage[most]{tcolorbox}
\usepackage{hyperref,fancyhdr,titlesec,capt-of,needspace}
\hypersetup{colorlinks=true,linkcolor=blue!50!black,urlcolor=blue!50!black}
\renewcommand{\contentsname}{Contenido}
\renewcommand{\figurename}{Figura}
\renewcommand{\tablename}{Tabla}
\sisetup{output-decimal-marker={.}}
\definecolor{ansys}{RGB}{255,183,27}
\titleformat{\section}{\Large\bfseries\color{black}}{\colorbox{ansys}{\makebox[1.6em]{\thesection}}}{0.6em}{}
\pagestyle{fancy}\fancyhf{}
\lhead{\small Placa con agujero a tensión --- Instructivo Ansys Mechanical}
\rhead{\small Ing. Juan Diego Rubio Ruiz}
\cfoot{\thepage}
\newtcolorbox{nota}[1][Nota]{colback=yellow!8,colframe=ansys!90!black,fonttitle=\bfseries,title=#1}
\newtcolorbox{alerta}[1][Cuidado]{colback=red!4,colframe=red!60!black,fonttitle=\bfseries,title=#1}
\newcommand{\menu}[1]{\textsf{\textbf{#1}}}
\newcommand{\paso}[2]{\par\needspace{0.42\textheight}\begin{center}\includegraphics[width=\linewidth,height=0.36\textheight,keepaspectratio]{img/#1}\captionof{figure}{#2}\end{center}}
\setlist[enumerate]{itemsep=2pt,topsep=4pt}

\begin{document}
\begin{titlepage}
\centering\vspace*{2cm}
{\Huge\bfseries Placa con agujero a tensión\par}\vspace{0.4cm}
{\LARGE Instructivo paso a paso en Ansys Mechanical 2026 R1 Student\par}\vspace{1cm}
\includegraphics[width=0.85\linewidth]{img/p19_resultado_malla5.png}\par\vspace{1cm}
{\large Análisis estático 2D en esfuerzo plano, estudio de convergencia de malla\\ y extracción del perfil de esfuerzos para validación analítica\par}
\vfill
{\large Ing. Juan Diego Rubio Ruiz\par}
{\normalsize Proyecto 1 del portafolio de simulación --- Octubre 2026\par}
\end{titlepage}

\tableofcontents
\newpage

\section{El problema}
Placa de acero estructural con un agujero circular central sometida a tensión uniaxial. Se busca el esfuerzo normal máximo $\sigma_{x,\max}$ en el borde del agujero.

\begin{center}
\begin{tabular}{ll}\toprule
Parámetro & Valor\\\midrule
Largo $L$ / ancho $W$ / espesor $t$ & \SI{200}{mm} / \SI{100}{mm} / \SI{5}{mm}\\
Diámetro del agujero $d$ & \SI{20}{mm} ($d/W=0.2$)\\
Material & Structural Steel ($E=\SI{200}{GPa}$, $\nu=0.3$)\\
Carga & $\sigma=\SI{10}{MPa}$ de tensión en los extremos\\
Modelo & 1/4 de placa (doble simetría), 2D esfuerzo plano\\
Referencia & Heywood: $\sigma_{\max}\approx\SI{31.4}{MPa}$\\\bottomrule
\end{tabular}
\end{center}

\begin{nota}[Organización de archivos]
Todo vive en \texttt{Documentos\textbackslash Proyectos\_Ansys\_Portafolio\textbackslash 01\_Placa\_con\_agujero\_FEA\textbackslash} con las subcarpetas \texttt{workbench}, \texttt{resultados}, \texttt{resultados\textbackslash imagenes}, \texttt{validacion} e \texttt{informe}.
\end{nota}

\section{Workbench: sistema de análisis en 2D}
\paso{p01_workbench_2d.png}{Crear el sistema Static Structural y fijar el tipo de análisis en 2D.}
\begin{enumerate}
\item Abrir Workbench 2026 R1 y arrastrar \menu{Static Structural} desde el \menu{Toolbox} al área central.
\item Guardar de inmediato: \menu{File $\rightarrow$ Save As} en la carpeta \texttt{workbench} con el nombre \texttt{placa\_agujero}.
\item Dar \textbf{un solo clic} en la celda \menu{Geometry}. Con \menu{View $\rightarrow$ Properties} activo, en \menu{Advanced Geometry Options} poner \menu{Analysis Type = 2D}.
\end{enumerate}
\begin{alerta}
El tipo de análisis 2D se fija \textbf{antes} de abrir la geometría o Mechanical. Después ya no se puede cambiar.
\end{alerta}

\section{Geometría en SpaceClaim}
\subsection{Plano de boceto XY}
\paso{p02_spaceclaim_boceto.png}{Modo boceto en SpaceClaim: elegir el plano XY y verlo de frente.}
\begin{enumerate}
\item Clic derecho en \menu{Geometry $\rightarrow$ New SpaceClaim Geometry}.
\item En la pestaña \menu{Boceto}, el segundo icono de la barra inferior es \menu{Seleccionar nuevo plano de boceto}. Dar clic en el eje Z: el plano queda perpendicular a Z, es decir, el plano XY.
\item \menu{Vista de plan}. El triedro debe mostrar X a la derecha y Y hacia arriba.
\end{enumerate}
\begin{alerta}
Mechanical exige que un modelo 2D esté en el plano XY. Si el triedro muestra la letra Z en el dibujo, el boceto está en otro plano.
\end{alerta}

\subsection{Rectángulo y círculo}
\paso{p03_boceto.png}{Rectángulo de $100\times50$ y círculo de diámetro 20 con centro en el origen.}
\begin{enumerate}
\item \menu{Rectángulo}: clic en el origen, mover el ratón hacia arriba a la derecha \textbf{sin dar otro clic}, teclear \texttt{100}, \texttt{Tab}, \texttt{50}, \texttt{Enter}.
\item \menu{Círculo}: clic en el origen, teclear \texttt{20} (diámetro), \texttt{Enter}.
\item \menu{Terminar edición del boceto} (palomita verde).
\end{enumerate}

\subsection{Quitar el círculo}
\paso{p04_borrar.png}{Borrar las caras del círculo con la herramienta Seleccionar.}
\begin{enumerate}
\item Activar \menu{Seleccionar} (no \menu{Tirar}: arrastrar con Tirar le da espesor a la superficie).
\item Clic en la parte del círculo fuera de la placa $\rightarrow$ \texttt{Supr}. Clic en el cuartito dentro de la placa $\rightarrow$ \texttt{Supr}.
\end{enumerate}
\paso{p05_geometria_final.png}{Geometría final: un solo cuerpo de superficie.}
Guardar con \texttt{Ctrl+S} y cerrar SpaceClaim. En Workbench, la celda Geometry queda con palomita verde.

\section{Mechanical: modelo}
\subsection{Unidades y comportamiento 2D}
\paso{p06_mechanical_unidades.png}{Unidades métricas en mm y esfuerzo plano.}
\begin{enumerate}
\item Doble clic en la celda \menu{Model} de Workbench para abrir Mechanical.
\item \menu{Home $\rightarrow$ Units $\rightarrow$ Metric (mm, kg, N, s, mV, mA)}.
\item En \menu{Geometry}, \menu{2D Behavior = Plane Stress}. El área reportada, \SI{4921.6}{mm^2}, coincide con $100\cdot50-\pi\cdot10^2/4$.
\end{enumerate}

\subsection{Espesor y material}
\paso{p07_espesor.png}{Espesor del cuerpo de superficie.}
Desplegar \menu{Geometry}, seleccionar el cuerpo y en \menu{Details} poner \menu{Thickness = 5 mm}. Revisar que \menu{Material $\rightarrow$ Assignment} sea \menu{Structural Steel}.

\subsection{Condiciones de simetría y carga}
\paso{p08_condiciones.png}{Mapa de aristas: dos simetrías y la carga.}
\begin{center}
\begin{tabular}{lll}\toprule
Arista & Objeto & Valores\\\midrule
Izquierda ($x=0$) & Displacement & $X=0$, $Y=$ Free\\
Inferior ($y=0$) & Displacement 2 & $X=$ Free, $Y=0$\\
Derecha ($x=100$) & Pressure & Magnitude $=\SI{-10}{MPa}$\\
Arco del agujero & --- & libre\\\bottomrule
\end{tabular}
\end{center}
Activar antes el filtro de selección \menu{Edge} en la barra de selección.

\paso{p09_disp1.png}{Displacement en la arista izquierda.}
\begin{enumerate}
\item Clic derecho en \menu{Static Structural (A5) $\rightarrow$ Insert $\rightarrow$ Displacement}.
\item Seleccionar la arista izquierda $\rightarrow$ \menu{Apply}.
\item \menu{X Component = 0}, \menu{Y Component = Free}.
\end{enumerate}
\paso{p10_disp2.png}{Displacement 2 en la arista inferior.}
Igual que el anterior en la arista inferior, con \menu{X = Free} y \menu{Y = 0}.

\paso{p11_pressure.png}{Presión en la arista derecha.}
Clic derecho en \menu{Static Structural $\rightarrow$ Insert $\rightarrow$ Pressure}, arista derecha $\rightarrow$ \menu{Apply}, \menu{Define By = Normal To}.
\paso{p12_signo.png}{Signo de la presión.}
\begin{alerta}[El signo importa]
En Ansys una presión positiva empuja hacia la cara (compresión). Para tensión se escribe \textbf{\texttt{-10}}. La leyenda debe decir ``Pressure: -10. MPa''.
\end{alerta}

\section{Malla y primera solución}
\paso{p13_malla.png}{Orden cuadrático y Edge Sizing en el arco.}
\begin{enumerate}
\item Clic en \menu{Mesh} $\rightarrow$ \menu{Element Order = Quadratic}.
\item Clic derecho en \menu{Mesh $\rightarrow$ Insert $\rightarrow$ Sizing}. Seleccionar el arco $\rightarrow$ \menu{Apply}. \menu{Type = Number of Divisions}, \menu{Behavior = Hard}.
\item Clic derecho en \menu{Mesh $\rightarrow$ Generate Mesh}.
\end{enumerate}
\paso{p14_solve.png}{Estadísticas de malla, resultado y Solve.}
\paso{p15_normal_stress.png}{Esfuerzo normal en X.}
\begin{enumerate}
\item En \menu{Mesh $\rightarrow$ Statistics} se leen \menu{Nodes} y \menu{Elements}.
\item Clic derecho en \menu{Solution (A6) $\rightarrow$ Insert $\rightarrow$ Stress $\rightarrow$ Normal}, \menu{Orientation = X Axis}.
\item \menu{Solve}.
\end{enumerate}
\paso{p16_resultado_malla_gruesa.png}{Primer resultado con la malla por defecto.}
Se compara contra $\sigma_x$ y no contra von Mises: en el punto $(0,10)$ el esfuerzo radial es nulo, así que $\sigma_x$ es directamente el esfuerzo máximo y es el mismo componente que da la teoría.

\section{Estudio de convergencia de malla}
\begin{alerta}[Lección aprendida]
Refinar \textbf{solo} el arco no converge. Con el tamaño global por defecto (\SI{8.77}{mm}), pasar de 4 a 8 divisiones \emph{bajó} $\sigma_{x,\max}$ de 27.78 a \SI{26.93}{MPa}: los elementos junto al agujero siguen siendo largos y deformados y no capturan el gradiente. Hay que refinar el tamaño global junto con el arco.
\end{alerta}
\paso{p17_element_size.png}{Los dos controles que se refinan juntos.}
\paso{p18_ciclo.png}{Ciclo por malla: cambiar tamaños, Solve, leer resultados.}

\begin{center}
\begin{tabular}{cccrrc}\toprule
Malla & Element Size & Divisiones & Nodos & Elementos & $\sigma_{x,\max}$ [MPa]\\\midrule
1 & \SI{4}{mm} & 4 & 1\,035 & 320 & 31.007\\
2 & \SI{2}{mm} & 8 & 3\,829 & 1\,228 & 31.310\\
3 & \SI{1}{mm} & 16 & 15\,144 & 4\,955 & 31.453\\
4 & \SI{0.7}{mm} & 24 & 30\,838 & 10\,147 & 31.486\\
5 & \SI{0.5}{mm} & 32 & 59\,977 & 19\,806 & 31.523\\\bottomrule
\end{tabular}
\end{center}
Los resultados se anotan en
\path{resultados\convergencia_malla.csv} y se procesan con \path{validacion\referencia_kirsch.py}. La versión Student admite hasta 128\,000 nodos en estructural.
\paso{p19_resultado_malla5.png}{Resultado de la malla 5.}

\section{Exportar resultados}
\subsection{Imágenes}
\paso{p20_image_to_file.png}{Image to File Preferences.}
\begin{enumerate}
\item Seleccionar el resultado o la malla, encuadrar con \menu{Box Zoom} (o \texttt{F7} para la vista completa).
\item \menu{Home $\rightarrow$ Images $\rightarrow$ Image to File}. Dejar \menu{Current Graphics Display} marcado y \menu{Append Graph} sin marcar (salvo para el Path) $\rightarrow$ \menu{OK}.
\item Guardar como PNG en \path{resultados\imagenes}: \path{sigmax_malla5_zoom}, \path{sigmax_malla5_completa}, \path{malla5_zoom}.
\end{enumerate}

\subsection{Perfil de esfuerzo sobre la arista izquierda}
\paso{p21_path.png}{Path sobre la arista $x=0$.}
\begin{enumerate}
\item Clic derecho en \menu{Model (A4) $\rightarrow$ Insert $\rightarrow$ Construction Geometry}; luego clic derecho en ella $\rightarrow$ \menu{Insert $\rightarrow$ Path}.
\item \menu{Path Type = Edge}. Con el filtro de aristas, seleccionar la arista izquierda (\texttt{F7} si el zoom no la deja ver) $\rightarrow$ \menu{Apply}.
\end{enumerate}
\paso{p22_normal_stress_path.png}{Esfuerzo normal evaluado sobre el Path.}
\paso{p23_path_resultado.png}{Perfil $\sigma_x(y)$ y exportación a texto.}
\begin{enumerate}
\item Insertar otro \menu{Normal Stress} con \menu{Scoping Method = Path}, el Path creado y \menu{Orientation = X Axis}.
\item Clic derecho $\rightarrow$ \menu{Evaluate All Results}.
\item Clic derecho $\rightarrow$ \menu{Export $\rightarrow$ Export Text File} como \texttt{resultados\textbackslash path\_sigma\_x.txt}. Para la imagen con la curva, \menu{Image to File} con \menu{Append Graph} marcado.
\end{enumerate}
\begin{nota}
El Path empieza arriba: $S=0$ corresponde a $y=50$ y $S=40$ al borde del agujero. El script usa la columna de coordenada $Y$, así que el sentido no importa.
\end{nota}

\section{Cierre}
\texttt{Ctrl+S} en Mechanical, cerrar, \texttt{Ctrl+S} en Workbench. Luego ejecutar \path{python validacion\referencia_kirsch.py} para regenerar las gráficas de convergencia y de comparación contra Kirsch.

\begin{nota}[Resumen de errores comunes]
\begin{itemize}[leftmargin=*]
\item Boceto en el plano XZ en lugar de XY.
\item Analysis Type en 3D por no cambiarlo antes de abrir la geometría.
\item Presión positiva (compresión) en lugar de \texttt{-10}.
\item Refinar solo el arco y dejar el tamaño global por defecto.
\item Usar la herramienta Tirar en lugar de Seleccionar al borrar caras.
\end{itemize}
\end{nota}
\end{document}
